\section{Conclusion}
\label{sec:conclusion}

We have developed a generalized hierarchical Bayesian method for multiple-changepoint segmentation with irregular designs, multiple aligned sequences, and known or latent group structure. For regular exponential-family segment models with proper conjugate priors, sufficient statistics give exact segment marginal likelihoods and posterior moments. These quantities enter common dynamic-programming recursions for the posterior distribution over segment counts and changepoint locations. A separate max-sum recursion with backtracking returns the joint MAP partition, so marginal posterior summaries and joint optimization are not conflated.

The same partition formulation accommodates design-dependent segment cohesion and boundary hazards, exact pooling of conditionally independent sequences with common changepoints, and a finite latent-group criterion optimized by alternating minorization and maximization. Nonconjugate segment models can be incorporated through numerical marginal-likelihood calculations; the resulting posterior approximation is related to the intended posterior only when a segmentwise error bound such as Assumption~\ref{ass:block-error} is established. Posterior prediction is likewise defined by the fitted observation family and its coordinate support.

The reported numerical studies examine several observation families, Gaussian boundary calibration, a finite latent-group simulation, computational scaling, and four applied datasets. In the reported clustering configuration, hard allocation accuracy is 96\%; this quantity characterizes the specified simulation, not a general recovery guarantee. The timing measurements likewise describe the tested implementation and sequence lengths, not an empirical proof of the worst-case complexity bound. For the applied datasets, changepoint agreement without an independently validated reference is not treated as accuracy against ground truth.

The additional protocols in Section~\ref{sec:experimental-design} address limitations of the present empirical evidence, including comparator tuning, prior sensitivity, model misspecification, numerical integration accuracy, and uncertainty under replicated simulations. Until these studies are completed, no comparative or general robustness claim follows from their planned analyses. The statistical contribution remains the integration of conjugate segment models, exact posterior calculations for locally factorized partitions, irregular-coordinate priors, shared-boundary pooling, and a finite-template clustering extension under explicitly stated assumptions.
